\begin{result}{F.3.27}{All-genus normalizer recognition for a whole absolutely simple Jacobian over Q}{Proved; every prime and finite arithmetic base change}
Let $C/\Q$ be a smooth proper geometrically connected curve
of genus $g\geq3$, with absolutely simple Jacobian and
$y\in C(\Q)$. Put $U=C\setminus\{y\}$.
For every prime $\ell$ and finite extension $K/\Q$,
a nontrivial closed procyclic subgroup of
$P=\pi_1(U_{\overline K})^{(\ell)}$ whose arithmetic
normalizer has open image in $G_K$ is contained in
an actual cusp inertia subgroup.
\end{result}
\begin{proof}
We retain the original curve and its full homology while
constructing every conjugate normalizer and universal map.

\paragraph{The first Magnus term and an actual finite detector.}
Let $A$ be the subgroup and $N$ its arithmetic normalizer.
The group $P$ is free pro-$\ell$ of rank $2g$, hence
torsion-free, so $A\simeq\Z_\ell$. For a generator $u$
its rational Magnus logarithm has a first nonzero term
$u_d\in\mathbb L_d(V)$, where $V=V_\ell J_C$.
Its normalizer character $\theta$ is trivial on
$N\cap P$, because geometric conjugation fixes that
initial term. It therefore gives an arithmetic
character on the open image of $N$.

The exact finite-cover construction in F.3.21 works
in this rank: take a finite $\ell$-quotient where
$u$ is nontrivial, and the preimage of its cyclic
image. The resulting actual finite etale cover $Y$
has nonzero rational homology class $[u]$ after
finite arithmetic descent. The normalizer intersection
with its arithmetic covering group still has open
image. Its geometric kernel acts inner and trivially
on the covering homology. Thus $[u]$ and $u_d$
have the very same arithmetic character $\theta$.
The original curve has one puncture, so $V$ is
pure of weight one and $\mathbb L_d(V)$ of weight
$d$. A nonzero proper-homology image of $[u]$
would force $d=1$ by a common good Frobenius.
It remains to exclude this degree-one case on
the original $V$.

\paragraph{The full CM field and all local primes.}
For $d=1$, the nonzero line $\Q_\ell u_1\subset V$
is stable on an open arithmetic subgroup. After
defining every geometric endomorphism, Faltings'
semisimplicity and endomorphism theorem identify
its projector with an idempotent of
$\operatorname{End}^0_{\Qb}(J_C)\otimes\Q_\ell$.
The rank formula in F.3.18 applies to this very
Jacobian. Absolute simplicity forces its full
geometric endomorphism algebra to be a CM field
$E$ of degree $2g$; the resulting simple quotient
is isogenous to $J_C$ itself.

The general CM/Q part of F.3.25 makes $E/\Q$
Galois and the actual endomorphism action
$G_{\Q}\to\operatorname{Aut}(E)$ surjective.
The module $V$ is rank one over
$E\otimes\Q_\ell=\prod_{\lambda\mid\ell}E_\lambda$.
Every idempotent rank is a sum of the positive
local degrees. Those degrees are all equal by
Galois transitivity, even at ramified primes.
A rank-one projector therefore requires complete
splitting of $\ell$ in $E$. This already excludes
$d=1$ at every nonsplit prime, with no prime
excluded from the argument.

\paragraph{All actual normalizer lines at a split prime.}
At a split prime the initial line is one of
the $2g$ embedding lines. Choose a nonzero
rational tangent vector at $y$ and its actual
arithmetic section. The exact formula
\[
 \rho_\ell(\tau)V_\varphi
          =V_{\varphi\circ a_\tau^{-1}}
\]
from F.3.25 supplies $2g$ elements $\tau_i\in G_{\Q}$
whose section lifts carry that line onto every
embedding line. Conjugate $A$ and $N$ by those
lifts. All resulting groups $A_i,N_i$ lie in
the arithmetic group of the same $U/\Q$.
Choose a finite Galois extension $L/\Q$ containing
$K$ and a field defining all endomorphisms, with
$G_L$ contained in every conjugate open image
$p(N_i)$. Every restricted $N_i$ now surjects
onto one common $G_L$.

Rosati duality pairs exactly the conjugate
embedding lines and is nondegenerate on each
pair, so generators of these $2g$ actual orbit
lines give a symplectic basis after scaling.
The arithmetic quotient of Nakamura's original
\S4, pp.~115--117, applies to every hyperbolic
$(g,n)$ with $n\geq1$. In particular Lemma~4.5,
Corollary~4.6, the degree-six rank table, and
the actual maps (4.7.1)--(4.7.2) give, for this
$(g,1)$ and this finite base $L$, an actual
$\sigma\in G_L(6)$ whose canonical special
derivation $D\in U_{g,6}$ satisfies
\[
                            q_{g,6}(D)\ne0.
\]
The quotient has dimension one and character
$\chi_\ell^3$ in this genus range
\cite[pp.~115--117]{AuditNakamura1996}.

Lift this same $\sigma$ to each $N_i$.
Its action on $V$ is the identity, so its
normalizer exponent on each $A_i$ is one.
The actual canonical representative and the
normalizing representative differ by an inner
automorphism of the punctured group. The rational
completed free-Lie calculation (QCM.4)--(QCM.5)
in F.3.25, now repeated on these $2g$ logarithms,
removes the degree-one inner term by a rational
power of the full subgroup logarithm and then
removes degrees two through five by the
centralizer of its nonzero linear term. It gives
\[
                         D(v_i)\in[\mathbb L_6(V),v_i]
                         \qquad(1\leq i\leq2g).
\]
Thus the same actual $D$ lies in
$U_{g,6}\cap T_{g,\mathrm{all}}$. F.3.26 proves
that its quotient is zero, a contradiction.
The degree-one case is excluded also at every
split prime.

\paragraph{The boundary character and the full normalizer.}
Return to the actual detector $Y$. Its nonzero
line $[u]$ has zero image in proper homology
because that would force the excluded case
$d=1$. It belongs to the actual boundary kernel
\[
 \left(\bigoplus_{t\in(\overline Y\setminus Y)(\Qb)}
          \Q_\ell(1)e_t\right)
 \big/\Q_\ell(1)\Bigl(\sum_t e_t\Bigr).
\]
Over a finite extension defining those points,
this entire kernel has character $\chi_\ell$.
Hence $\theta=\chi_\ell$ on the intersected
normalizer. Any element of the full normalizer
over that extension has the same arithmetic
image as an element of the intersected one;
their quotient lies in $N\cap P$, where
$\theta$ is trivial. Thus the full normalizer
also acts cyclotomically and has open arithmetic
image. Nakamura's Theorem~2.1.1, p.~94, for
the full class of finite $\ell$-groups gives
actual cusp inertia containing $A$
\cite{AuditNakamura1994}. The finite arithmetic
extensions change neither the geometric group
nor this unique-puncture inertia conjugacy class.
\end{proof}

\paragraph{Exact relation to the general detecting-cover problem.}
The assertion concerns the whole absolutely simple
Jacobian of a specified once-punctured curve over
$\Q$. The rational normal-core construction of
F.3.23 gives an actual once-punctured $\Q$-curve,
but its Jacobian need not be absolutely simple.
For a proper CM quotient the conjugate detected
lines still lie in its actual images inside the
larger curve homology. The mixed normalizer-trace
equations or an actual punctured target map
needed there are retained as unresolved steps;
the full LM6 A.22 subgroup assertion is unchanged.
